\chapter{Um Momento Para Amar}

\lettrine{C}{om a ênfase que o budismo} coloca na reflexão, na atenção plena e
na sabedoria, a Vida Santa pode, por vezes, parecer uma tentativa quase
insensível de olhar para tudo de forma muito objetiva. Em vez de sentir as
coisas, devemos ver tudo como \emph{anicca, dukkha} e \emph{anatta.} É assim que
pode parecer. Mas lembre-se, a experiência sincera da vida é amorosa, pelo que
o amor e a devoção não devem ser descartados.

Se encararmos a experiência do amor apenas como \emph{anicca, dukkha} e
\emph{anatta}, isso pode parecer insensível. A objetividade, no entanto, é
apenas a forma de ver as coisas com perspetiva, para que o amor não seja algo
que nos cegue. Se estivermos apegados à ideia do amor, podemos ficar bastante
cegos à sua realidade. Podemos ficar muito inspirados ao falar sobre ele ou ao
meditar sobre o amor -- procurando-o nos outros, exigindo-o ou sentindo-nos,
de alguma forma, excluídos. Mas o que é o amor em termos das nossas vidas, tal
como as vivemos?

A nível emocional, pode querer sentir uma tremenda sensação de unidade, ou
talvez direcionar esses sentimentos para uma pessoa em particular, desejando ter
uma relação especial e amorosa com outra pessoa. O amor também pode ser abstrato
-- amor por todos os seres humanos, amor por todos os seres, amor por Deus, amor
por algo ou por algum conceito.

A devoção vem do coração, não é algo racional. Não se pode obrigar-se a sentir
amor ou devoção apenas porque se gosta da ideia. É quando não se está apegado,
quando o coração está aberto, recetivo e livre, que se começa a experimentar o
que é o amor puro. Amor benevolente, compaixão, alegria solidária, equanimidade
-- o reino das Moradas Divinas, os \emph{brahma-viharas} -- estes provêm de uma
mente vazia. Não de uma posição estéril de mero sentimento de aniquilação, mas
de um coração que não está iludido, não cegado por ideias do eu ou dos outros,
ou por paixões de algum tipo.

Podes pensar que a Vida Santa é fria e sem coração porque, numa comunidade de
\emph{samaṇas} como esta, vivendo de acordo com um caminho de contenção e
disciplina, não somos expressivos nas nossas manifestações de amor e alegria.
Esta comunidade não transborda de sentimentos de devoção. É bastante formal e
contida na sua forma e na sua expressão. Mas isso não nega necessariamente o
amor. Com a atenção plena, com a forma como nos relacionamos -- com os nossos
próprios corpos, com o Sangha, para com os leigos, para com a tradição e para
com a sociedade -- há uma abertura, bondade e recetividade. Há carinho, alegria
e compaixão que podemos sentir.

Continua a ser \emph{anicca} e \emph{anatta}, e é \emph{dukkha} no sentido de
que não é, em si mesmo, o fim de nada; não é satisfatório como identidade ou
apego. Mas quando o coração está livre das ilusões do eu, surge então uma
qualidade amorosa na pura alegria de ser. É apenas a maneira natural das coisas.
Portanto, quando contemplamos dessa forma, esse é o caminho da fé, da confiança
e da devoção.

Quando falamos de fé, confiança e segurança, não são coisas que se possa
realmente compreender. A fé não é algo que se possa criar. Pode-se dizer as
palavras, mas ter realmente fé e confiança no Dhamma é estar disposto a
abandonar qualquer exigência, afirmação ou qualquer tipo de apego. E essa
experiência de fé surge à medida que examinamos e compreendemos o Dhamma, ou a
verdadeira natureza das coisas. Se realmente contemplarmos o Dhamma, se o
virmos, então surge a fé, esse forte sentimento de confiança total, de segurança
na verdade.

\emph{Se} estás a praticar meditação \emph{Vipassana} e te estás a sentir mais
assustado, ansioso ou tenso, ou a sentir-te emocionalmente esvaziado, então não
o estás a fazer da maneira certa. Talvez estejas a usar a técnica como uma forma
de suprimir os teus sentimentos ou de negar as coisas. Assim, acabas por te
sentir mais tenso, cético, incerto. Existe um apego a uma certa visão sobre
isso. Quanto mais vemos e compreendemos completamente como as coisas são, mais
temos essa qualidade de fé. Essa fé aumenta, é uma confiança total. Quando se
fala em rendição, desistir ou deixar ir, é através da confiança total. Não se
trata apenas de arriscar, mas sim da experiência da fé.

O caminho é algo que cultivamos. Temos de saber onde estamos e não tentar
tornar-nos algo que pensamos que gostaríamos de ser; temos de praticar com a
forma como as coisas são agora, sem fazer um julgamento sobre isso. Se te sentes
tenso, nervoso, desiludido, desapontado contigo mesmo ou com a tradição, ou com
o professor ou com os monges ou monjas, ou o que quer que seja, então tenta
reconhecer que o que está no momento é suficiente. Esteja disposto a
simplesmente admitir, a reconhecer as coisas como elas são, em vez de se
entregar à crença de que o que está a sentir é, de alguma forma, uma descrição
precisa da realidade ou que está errado e não deveria estar a sentir-se assim.
Esses são dois extremos. Mas o cultivo do Caminho consiste em reconhecer que
tudo o que está sujeito a surgir está sujeito a cessar. E isto não é uma forma
depreciativa ou fria de cultivar o caminho, mesmo que possa parecer isso.

Podem pensar que basta libertares-te de todos os teus sentimentos e perceberes
que o amor no teu coração é \emph{anicca, dukkha, anatta.} Sentes amor pelo Buda
e pensas: «Oh, isso é apenas \emph{anicca, dukkha, anatta.} É só isso!» Sentes
amor pelo professor e pensas: «Isso é apenas \emph{anicca, dukkha, anatta.} Não
te apegues ao professor!» Sentes amor pela tradição\ldots{}
``\emph{anicca, dukkha, anatta,} não te apegues às tradições -- ou às técnicas.''

«Não se apegar a nada» pode ser apenas uma forma de suprimir \emph{tudo.} Não é
necessariamente desapego ou não-apego, pode ser apenas uma posição que se
assume. E se assumires essa posição e agires a partir dela, tudo o que vais
sentir é negatividade e stress. «Não deves apegar-te a nada; não deves amar
nada, não deves sentir nada -- sentir qualquer coisa é apenas \emph{anicca,
  dukkha, anatta}.» Isso significa que estás apenas a pegar nas palavras e a
usá-las como um porrete, um grande bastão contra a tua mente. Não estás a
refletir, a abrir-te e a confiar.

A prática de \emph{Mettā} é uma das belas práticas devocionais altamente
recomendadas no Buddha-Dhamma. Bondade amorosa. Como seres humanos, somos
criaturas de sangue quente. Nós sentimos amor. Isso faz parte da nossa
humanidade. Gostamos uns dos outros; gostamos de estar com pessoas; gostamos de
ser gentis; sentimos prazer em cozinhar e oferecer comida a outras pessoas.
Gostamos de ajudar. Isso é visível no costume do \emph{dana} nas comunidades
asiáticas. Quando os cingaleses vêm cá com os seus caris, ficam radiantes. É a
alegria de dar.

Ora, isso é uma qualidade muito boa, não é? É lindo ver alguém que pode ter
passado a noite em claro a preparar comida deliciosa para oferecer a outra
pessoa -- não a está a cozinhar para si próprio. Bem, o que é isso enquanto
experiência humana? É impureza ou é estar apegado ao prazer ou à felicidade de
fazer coisas pelos outros? Esta é a beleza da humanidade, não é? -- simplesmente
ser capaz de amar, de dar, de partilhar, de ser generoso.

Tente refletir sobre qual seria o grande prazer de ser a pessoa mais rica do
mundo? O que seria verdadeiramente gratificante? Conseguir o que quero? Não,
seria a oportunidade de o partilhar, não seria? Esse seria o verdadeiro prazer
de ser rico e abastado -- poder partilhá-lo, como \emph{dana}, generosidade; ao
passo que ser rico e não poder partilhá-lo seria um verdadeiro fardo. Que fardo
seria esse, ser o homem mais rico do mundo e ser egoísta, agarrando-me a tudo e
guardando tudo só para mim!

A alegria da riqueza está na capacidade de a partilhar e de a dar sem qualquer
tipo de intenções corruptas ou exigências egoístas.

Então, isto é o que há de encantador na nossa humanidade: podemos experimentar
esta alegria de dar. E é algo que todos experimentamos quando realmente damos
algo, quando ajudamos alguém sem qualquer pedido egoísta ou exigência de algo em
troca. Então, experimentamos alegria. É certamente uma experiência humana
encantadora -- mas não esperamos que ela nos torne felizes para o resto das
nossas vidas.

A alegria da generosidade e da bondade não é permanente, não nos torna
permanentemente felizes -- mas não esperamos que o faça. Se o fizéssemos, já não
seria \emph{dana}, seria um acordo que estávamos a fazer. Não seria um ato de
generosidade, seria comprar algo.

A verdadeira alegria vem de dar e não se importar se alguém sabe ou reconhece
isso. Assim que o ego entra, por exemplo: «Estou a dar este \emph{dana} -- EU, estou a dar!»,
então a quantidade de alegria que vem de dar é provavelmente
muito mínima. Se estou tão preocupado que reconheças e aprecies a minha
generosidade e a minha bondade, então isso torna-se um sabor de mente sem
alegria. Não se pode sentir-se feliz ou ter verdadeira alegria na vida se
houver apego à ideia de que as nossas ações devem ser reconhecidas. Não há nada
de errado em as pessoas apreciarem a bondade e a generosidade de outra pessoa --
mas quando não a exigimos, então há alegria.

O amor romântico baseia-se geralmente na ilusão de um eu e na exigência de algo
em troca. O amor espiritual, então, é amor altruísta ou amor universal e é
representado pelos \emph{brahma-viharas -- mettā, karuṇā, muditā} e
\emph{upekkhā.} Esse amor é uma experiência unificadora. Ele aproxima, une. É
uma comunhão. O ódio é a experiência da separação. Quando odiamos, não há união,
comunhão ou unidade. O ódio é separador, divisivo e discriminatório. O amor é
unitivo e queremos unidade porque viver num mundo de ódio, discriminação e
separação é um reino infernal miserável.

A comunidade é uma comunhão, uma Sangha, um todo. Se estivermos separados da
Sangha, se odiarmos a Sangha, «Odeio esta monja, odeio aquele monge -- e não
gosto daquilo, não gosto disto», então isto não é comunidade, é desunião. Esse
sentimento de alienação, separação, que enfatiza o eu e o tu, as \emph{tuas}
falhas e os \emph{meus} sentimentos, e a minha raiva pelas tuas falhas. Também
posso ser eu a enfatizar as coisas que estão erradas em ti -- coisas que estão
erradas nos monges, coisas que estão erradas nas monjas, coisas que estão
erradas nos anagairikas, coisas que estão erradas, ponto final. E apegar-me a
essas perceções far-me-á sentir alienado, separado, zangado, descontente e
deprimido.

Às vezes, a mente entra num estado muito negativo em que tudo o que sentimos é
irritação. O que quer que as pessoas façam não parece suficientemente bom.
Quando se está nesse estado de espírito, tudo parece errado -- os gatos, o sol,
a lua -- a mente entra em divisão, separação e negatividade. Sentes-te separado
de tudo o que vês, e nenhuma comunhão ou união é possível enquanto estiveres
identificado e apegado a essa atitude mental. Quando estás num estado de
espírito amoroso, então não importa realmente se alguém não se está a sentir
muito bem ou se não está a fazer exatamente o que deveria. Há sempre pequenas
coisas, pequenos detalhes que não são bem o que deveriam ser. Mas quando estás
num estado de espírito amoroso, essas coisas não são tão importantes.

Assim, a experiência amorosa surge porque estás disposto a ignorar as diferenças
de personalidade e a discriminação que existe no reino condicionado, em prol do
sentimento de comunhão, de união, de unidade. Estamos a unir-nos como irmãos e
irmãs numa experiência comum de velhice, doença e morte, em vez de apontar as
diferenças ou quem é melhor do que quem. Quando nos refugiamos na Sangha,
estamos a refugiar-nos em \emph{supaṭipanno} (aqueles que praticaram bem),
\emph{ujupaṭipanno} (aqueles que praticaram diretamente), \emph{ñāyapaṭipanno}
(aqueles que praticaram com discernimento),
\emph{sāmīcipaṭipanno}\footnote{Estes termos fazem parte dos cânticos diários em
  homenagem ao Buddha, Dhamma e Sangha.} (aqueles que praticam com integridade).
Em vez de nos refugiarmos em americanos, britânicos, australianos, em homens ou
mulheres, em monjas ou monges, refugiamo-nos naqueles que praticam o Dhamma --
nos bons, nos diretos, nos sinceros.

Temos tendências tanto para a união como para a separação e podemos estar
conscientemente atentos a estas. A forma como as coisas são tem de ser
reconhecida como Dhamma. Há união e há separação, e com uma consciência clara
não nos identificamos com nenhum dos extremos. Há tempo para a união e a
comunhão; tempo para a não discriminação, para a devoção, para a gratidão, para
a generosidade, para a alegria.

Mas há também tempo para a separação e a discriminação, para examinar o que está
errado. É necessário olhar para as falhas -- olhar para a raiva, a inveja e o
medo -- e aceitar e compreender essas experiências emocionais, em vez de as
julgar e considerá-las como parte de si mesmo e como algo que não se deveria
ter.

É isto que significa ser humano: nascemos numa forma separada e, no entanto,
podemos unir-nos. Podemos alcançar a unidade, a comunidade e a unicidade; mas
também podemos discriminar.

Assim, o refúgio no Buda é a capacidade de um ser humano de reconhecer ambos os
lados e de responder de forma adequada. Podemos olhar para as falhas e os
problemas da vida como parte da nossa experiência humana, em vez de o fazermos
de forma pessoal. Já não estamos a proliferar, nem a ampliar ou a exaltar,
obcecados com o que está errado, porque temos esta perspetiva de unidade e
separação. É assim que as coisas são: o Dhamma.

Será que ser um monge budista ou uma monja budista é uma negação do amor? Será
que a disciplina do Vinaya é apenas um meio de suprimir sentimentos? \emph{Pode}
ser apenas isso. \emph{Podemos} usar a disciplina do Vinaya e a tradição
monástica apenas como uma forma de evitar as coisas. Talvez os monges tenham
apenas medo das mulheres\ldots{} Talvez as monjas estejam simplesmente
aterrorizadas com os homens, por isso tornam-se monjas e não têm de enfrentar
os seus medos e ansiedades no que diz respeito às relações com os homens\ldots{}

E, claro, muitas pessoas mundanas pensam assim, não é verdade? Acham que estamos
todos aqui por não conseguirmos lidar com o mundo real.

Mas será que é mesmo assim? Se for, se é por isso que se é monge ou monja, então
está-se nisso pelas razões erradas. Isto não é um instrumento para evitar a
realidade e a vida, mas para refletir sobre elas. Porque na contenção e na
dignidade da contenção, o caminho do monaquismo é uma expressão de amor por
\emph{todos} os seres -- homens, mulheres, tanto dentro como fora. Agora já não
escolhemos uma pessoa para concentrar a nossa atenção e a quem nos dedicamos,
mas dedicamo-nos a todos os seres.

Percebo que, se fosse um homem de família, toda a minha atenção teria de estar
voltada para a minha esposa, os meus filhos e a família imediata. Esse é o
resultado da vida familiar e o que o casamento implica. Eles têm prioridade.
Temos de nos relacionar com aqueles com quem estamos casados e pelos quais somos
responsáveis.

Pode-se ser um mendicante, viver apenas da fé, da confiança na bondade e
benevolência dos outros seres, porque se sente amor ou respeito por todos os
seres. O amor e o respeito por todos os seres é o que gera as esmolas que nos
sustentam nesta vida como mendicantes. E o engraçado é que o poder da Sangha
budista é tão forte que, mesmo que você pessoalmente odeie todos os outros
seres, as esmolas continuam a chegar! O poder do manto parece ser tão forte que,
mesmo que você, como indivíduo monge ou monja odeie toda a gente, continuará a
ser alimentado por seres de bom coração. Isto deve-se às \emph{paramitas} do
Buda. Isto não significa que deva desenvolver ódio ou justificá-lo em si mesmo
de forma alguma. Pelo contrário, é uma reflexão sobre o poder de uma convenção
muito hábil que foi estabelecida pelo Senhor Buda. Quando aprecia isso, então
sente realmente amor e confiança.

Por que é que estes mosteiros aqui em Inglaterra funcionam? Por que é que
deveriam funcionar num país não budista? Por que é que alguém iria querer enviar
um cheque pelo correio, trazer um saco de batatas ou preparar uma refeição? Por
que é que se dariam ao trabalho? Isto deve-se às \emph{paramitas} do Senhor
Buda. A bondade do estilo de vida que ele estabeleceu gera generosidade. A
bondade amorosa, a compaixão e a alegria da Vida Santa alcançam e abrem outras
pessoas para essa mesma experiência.

É um mistério. De uma atitude prática e mundana de justificar a nossa existência
aos olhos da sociedade, não parece que façamos assim tanto por ninguém. Muitas
pessoas pensam que ficamos aqui sentados a tentar alcançar a iluminação para nós
próprios -- a ter estados mentais agradáveis e prazerosos porque não suportamos
o mundo real. Mas quanto mais contemplamos esta vida e a compreendemos, mais
percebemos o poder da bondade, da fé, das \emph{paramitas} do Buda, que permitem
que uma comunhão tenha lugar numa interconexão de bondade.

E isso não precisa de ser demonstrado, discutido ou enfatizado em demasia. Fala
por si mesmo. Não precisamos de sair por aí a dizer às pessoas: «Devem dar-nos
esmolas porque estamos a praticar o Dhamma e somos discípulos do Buda.» As
nossas necessidades são supridas porque as pessoas apreciam e respeitam a Vida
Santa. Ela traz alegria e felicidade à vida das pessoas porque nos regozijamos
com a beleza dos outros e com a bondade e benevolência desta experiência de
viver. Na verdade, a Vida Santa é estranha, uma forma estranha de viver. A forma
como funciona é um mistério em termos do que consideramos realidade de acordo
com o nosso condicionamento cultural. Mas como Dhamma, como Verdade, como a
forma como as coisas são, funciona de facto. E isto aumenta a nossa fé e a nossa
confiança nos Refúgios e na beleza e bondade das nossas vidas como
\emph{samaṇas}.

\photoPage{image-016-stupa-in-snow.jpg}{Stupa na neve, Amaravati}{}
